\documentclass[10pt,a4paper]{amsart}
\usepackage[margin=19mm]{geometry}
\usepackage{amsmath,amssymb,mathtools}
\usepackage[scr=rsfs,cal=euler]{mathalfa}
\usepackage{xfrac}
\usepackage[unicode,hidelinks,bookmarks=false]{hyperref}
\newcommand{\ie}{i.e.\ }
\newcommand{\cf}{cf.\ }
\newcommand{\resp}{respectively\ }
\newcommand{\Subsection}{\subsection}
\providecommand{\nobreakdash}{\nobreak\hbox{-}}
\setlength{\emergencystretch}{2em}
\multlinegap0pt
\newtheorem{theorem}{Theorem}
\newtheorem{lemma}[theorem]{Lemma}
\newtheorem{corollary}[theorem]{Corollary}
\def\C{\mathbb C}
\def\P{\mathbb P}
\def\l{\mathcal}
\def\del{\partial}
\def\GL{\textup{GL}}
\def\Hom{\textup{Hom}}
\def\lp{\left(}
\def\rp{\right)}
\title[Lemma 2.13]{The regular singular inverse problem in differential Galois theory: the specialization lemma for regular singular equations (Lemma 2.13)}
\author{Thomas Serafini, Michael Wibmer}
\date{Source: arXiv:2501.07993v2 (CC BY 4.0)}
\begin{document}
\maketitle

\noindent\emph{Source and licence.} The regular singular inverse problem in differential Galois theory, by Thomas Serafini, Michael Wibmer. Version arXiv:2501.07993v2
(CC BY 4.0), distributed by arXiv under the Creative Commons Attribution 4.0 International licence,
http://creativecommons.org/licenses/by/4.0/, as recorded in the arXiv licence metadata for this exact
version and shown on the abstract page https://arxiv.org/abs/2501.07993v2. Full licence notice:
\texttt{licenses/arxiv-2501.07993-CC-BY-4.0.txt}.\par\smallskip

\noindent\textit{Editorial scope.} Counted unit: the article's Lemma 2.13 (source0.tex lines 666--668, printed as Lemma 2.13): for algebraically closed fields $k\subseteq k'$, a subset $S=\{p_1,\ldots,p_{d+1}\}$ of $\P^1(k)$ with $d+1$ elements and a matrix $A\in k'(z)^{n\times n}$ whose equation $\del(y)=Ay$ has all poles contained in $S$, is Fuchsian at $p_1,\ldots,p_d$ and regular singular at $p_{d+1}$, there exist a finitely generated $k$-subalgebra $\l B\subseteq k'$ and a monic $f\in\l B[z]$ with $A\in\l B[z]_f^{n\times n}$ such that every $k$-specialization $c$ of $\l B$ preserves all three properties. The article's complete original proof of it is delivered with it (source0.tex lines 670--743, one \texttt{proof} environment of 74 lines, adjacent to the statement, with no gap and no deferral). No mathematics is written, repaired or re-derived: the statement and the proof are the article's own lines, in the article's own notation, and no retained line is substituted, re-flowed or omitted, so nothing is removed and the package carries no omission record. Retained uncounted as the source-local context the proof needs, each exactly the statement of the result it invokes and nothing else: the local specialization lemma at a point (lines 629--631, printed as Lemma 2.12, source label \texttt{lem: specialize regsing at 0}) and the gauge-normal-form lemma for a regular singular equation at $0$ (lines 389--392, printed as Lemma 2.7, source label \texttt{lem:basis}). Every cross-reference of every retained line resolves inside this document, and no retained line contains a citation, so the wrapper carries no bibliography and no bibliography entry. No retained line contains a figure environment, an image-inclusion command or drawing code, so no image is delivered and the package declares no asset dependency.

Editorial formatting only, in addition: an AMS \texttt{amsart} preamble that restates the article's own declarations and macros used by the retained lines. The article's theorem-like environments and its mathematical macros are declared in the authors' own local style file \texttt{TSeng.sty}, which is shipped in the e-print but is not redistributed here; the wrapper restates only what the retained lines need, namely the environments \texttt{theorem}, \texttt{lemma} and \texttt{corollary} on one shared counter as in the article but without the article's section-based prefix, and the eight macros \texttt{\textbackslash C}, \texttt{\textbackslash P}, \texttt{\textbackslash l}, \texttt{\textbackslash del}, \texttt{\textbackslash GL}, \texttt{\textbackslash Hom}, \texttt{\textbackslash lp} and \texttt{\textbackslash rp}, copied verbatim from that file. The retained environments therefore print in this document as Lemma 1, Lemma 2 and Lemma 3 in order of appearance, whereas the article prints the same three items as Lemma 2.7, Lemma 2.12 and Lemma 2.13.

The complete native source of the article as distributed in the arXiv e-print for this exact version is bundled unchanged as \texttt{sources/171589/source0.tex}: it is byte identical to the e-print file \texttt{regsing.tex} except that the pipeline's mechanical concatenation prepends one header comment line naming it. The article's bibliography file \texttt{diffgaloisrefs.bib} and its local style file \texttt{TSeng.sty} are not redistributed.
\begin{lemma}\label{lem:basis}
	Let $A \in k(t)^{n\times n}\subseteq k((t))^{n\times n}$ be such that $\del(y)=Ay$ is regular singular at $0$. Then there
	exists a matrix $F\in \GL_n(k(t))$ such that $A'=FAF^{-1}+\del(F)F^{-1}$ satisfies $tA'\in k[t]_{(t)}^{n\times n}$. In other words, if $A$ is Gauge equivalent over $k((t))$ to a matrix $A'$ satisfying $tA'\in  k[[t]]^{n\times n}$, then $A$ is already Gauge equivalent over $k(t)$ to such a matrix.
\end{lemma}

\begin{lemma}\label{lem: specialize regsing at 0}
	Let $k\subseteq k'$ be an inclusion of algebraically closed fields and let $A \in k'(t)^{n \times n}\subseteq k'((t))^{n\times n}$ be such that $\del(y)=Ay$ is regular singular at $0$. Then there exists a finitely generated $k$\=/subalgebra $\l B$ of $k'$ and a monic polynomial $f \in \l B[t]$ such that $A \in \l B[t]_f^{n \times n}$ and $\del(y)=A^c y$ is regular singular at $0$ for every $c \in \Hom_k(\l B, k)$.	
\end{lemma}

\begin{lemma}\label{lem:regsing}	
	Let $k\subseteq k'$ be an inclusion of algebraically closed fields and let $S=\{p_1,\ldots,p_{d+1}\}$ be a subset of $\P^1(k)\subseteq\P^1(k')$ with $d+1$ elements. Furthermore, let $A \in k'(z)^{n \times n}$ be such that $\del(y)=Ay$ has all poles contained in $S$, is Fuchsian at $p_1,\ldots,p_d$ and regular singular at $p_{d+1}$. Then there exists a finitely generated $k$-subalgebra $\l B \subseteq k'$ and a monic polynomial $f \in \l B[z]$ such that $A \in \l B[z]_f^{n \times n}$ and for every $c \in \Hom_k(\l B, k)$ the differential equation $\del (y) =A^c y$ has all poles contained in $S$, is Fuchsian at $p_1,\ldots,p_d$ and regular singular at $p_{d+1}$.	
\end{lemma}

\begin{proof}
%	There is no difficulty in finding a finitely generated $k$-subalgebra $\l B_0$ of $k'$ and a monic $f_0\in\l B[z]$ such that $A \in \l B_0[z]_f^{n \times n}$ (cf. the proof of Lemma \ref{lem: specialize regsing at 0}).
	We first assume that $\infty\notin S$. By assumption,
	$B=(z-p_1)\ldots(z_1-p_d)(z-p_{d+1})^eA\in k'[z]^{n\times n}$ for some integer $e\geq 1$.	
%	As the poles of $\del(y)=Ay$ are contained in $S$, there exist (not necessarily distinct) $a_1,\ldots,a_m\in S$ such that $B=(z-a_1)\ldots(z_1-a_m)A$ has entries in $k'[z]$.
		 Let $\l B_0$ be the $k$\=/subalgebra of $k'$ generated by the coefficients of all entries of $B$ and set $f=(z-p_1)\ldots(z_1-p_d)(z-p_{d+1})^e\in k[z]\subseteq\l B_0[z]$.
	Then $A\in \l B_0[z]_{f}^{n\times n}$ and for any $c \in \Hom_k(\l B_0, k)$ the poles in $k$ of the differential equation $\del (y) =A^c y$ are contained in $S$.
	
	For a non-zero rational function $h\in \l B_0[z]_{f}\subseteq k'(z)$, when written as $h=\frac{h_1}{f^m}$, with $h_1\in\l B_0[z]$, the rational function $-\frac{1}{t^2}h(\frac{1}{t})\in k'(t)$ has a pole at $t=0$ if and only if $\deg(h_1)-m(d+e)+2>0$. Thus, if $-\frac{1}{t^2}h(\frac{1}{t})$ does not have a pole at $t=0$, for any specialization $c\in \Hom_k(\l B_0, k)$, the rational function $-\frac{1}{t^2}h^c(\frac{1}{t})$ does not have a pole at $t=0$. So, as $p=\infty$ is not a pole of $\del(y)=Ay$, 
	$p=\infty$ is also not a pole of $\del(y)=A^cy$ for any $c\in\Hom_k(\l B_0,k)$. In summary, for any $c\in\Hom_k(\l B_0,k)$, the poles of $\del(y)=A^cy$ are contained in $S$.
	
%	
%	For $f_0=(z-a_1)\ldots(z_1-a_m)$ we then have $A\in \l B_0[x]_{f_0}^{n\times n}$.
%	
%	
%	The singularities of $\del(y)=Ay$ are contained in $\{a_1,\ldots,a_m\}\cup\{\infty\}$. Without loss of generality we may assume that $a_1,\ldots,a_r$ are regular points for $\del(y)=Ay$ and that $a_{r+1},\ldots,a_{m}$ are singular points for $\del(y)=Ay$.
%	
	 The differential equation $\del(y)=A(t+p_{d+1})y$ is regular singular. Thus, by Lemma \ref{lem: specialize regsing at 0}, there exists a finitely generated $k$-subalgebra $\l B$ of $k'$ and a monic polynomial $f_1\in \l B[t]$ such that $A(t+p_{d+1})\in \l B[t]_{f_1}^{n\times n}$ and $\del(y)=A(t+p_{d+1})^cy$ is regular singular for all $c\in\Hom_k(\l B,k)$. Enlarging $\l B$ if necessary, we may assume that $\l B_0\subseteq \l B$. Furthermore, replacing $f_1\in \l B[t]$ by $f_1(t)f(t+p_{d+1})\in \l B[t]$, we may assume that $f(t+p_{d+1})$ divides $f_1$ in $\l B[t]$.
	 We claim that $\l B$ and $f\in\l B[z]$ have the required properties. 
	 
	 Let $c\in\Hom_k(\l B,k)$. Note that $\l B_0[t+p_{d+1}]_{f(t+p_{d+1})}\subseteq \l B[t]_{f_1}$ and $h(t+p_{d+1})^c=h^c(t+p_{d+1})$ for $h\in \l B_0[z]_{f}$. So  $A(t+p_{d+1})^c=A^c(t+p_{d+1})$.
%	 	 $$\l B_0[z]_{f_0}\subseteq \l B[z]_{f_0}=\l B[z-p_{d+1}]_{f_0}\subseteq \l B[z-p_{d+1}]_{f_1(z-p_{d+1})}=\l B[z]_f$$ and $A(t+p_{d+1})^c=A^c(t+p_{d+1})$ for every $c\in\Hom_k(\l B,k)$.
	  As $\del(y)=A(t+p_{d+1})^cy$ is regular singular, this shows that $\del (y)=A^cy$ is a regular singular at $p_{d+1}$. As explained in the second paragraph of this proof, the poles of $\del (y)=A^cy$ are contained in $S$.
	 Finally, from $(z-p_1)\ldots(z_1-p_d)(z-p_{d+1})^eA\in \l B[z]$, we obtain $(z-p_1)\ldots(z_1-p_d)(z-p_{d+1})^eA^c\in k[z]$, showing that $\del (y)=A^cy$ is Fuchsian at $p_1,\ldots,p_{d}$.
	 
	 In case $\infty\in S$, a slight modification of the above proof yields the result. In detail, if $\infty$ is among the Fuchsian points, say $p_1=\infty$, then, with $f=(z-p_2)\ldots(z-p_d)(z-p_{d+1})^e\in k[z]$, we can proceed as above, noting that $-\frac{1}{t}A(\frac{1}{t})\in k'[[t]]^{n\times n}$ implies $-\frac{1}{t}A^c(\frac{1}{t})\in k[[t]]^{n\times n}$ for any $c\in\Hom_k(\l B_0,k)$ (cf. the second paragraph of this proof).
	 
	 If $\infty$ is only regular singular, i.e., $p_{d+1}=\infty$, then, with $f=(z-p_1)\ldots(z-p_d)\in k[z]$ we can proceed similarly to the above, but using that, as $\del(y)=-\frac{1}{t^2}A(\frac{1}{t})y$ is regular singular, there exists, by Lemma~\ref{lem: specialize regsing at 0}, a finitely generated $k$-subalgebra $\l B$ of $k'$ and a monic polynomial $f_1\in\l B[t]$ such that $-\frac{1}{t^2}A(\frac{1}{t})\in \l B[t]_{f_1}^{n\times n}$ and $\del(y)=(-\frac{1}{t^2}A(\frac{1}{t}))^cy$ is regular singular for every $c\in\Hom_k(\l B,k)$. Enlarging $\l B$ if necessary, we can assume that $\l B_0\subseteq \l B$.  Writing $f(\frac{1}{t})=\frac{cf_2}{t^m}$ with $c\in k$ non-zero and $f_2\in k[t]$ monic and replacing $f_1$ by $f_1tf_2$, we can assume that $t$ and $f_2$ divide $f_1$ in $\l B[t]$. As $\l B_0[\frac{1}{t}]\subseteq \l B_0[t]_t\subseteq \l B[t]_{f_1}$ and $\frac{1}{f(\frac{1}{t})}=\frac{c^{-1}t^m}{f_2}\in \l B[t]_{f_1}$, we see that $\l B_0[\frac{1}{t}]_{f(\frac{1}{t})}\subseteq \l B[t]_{f_1}$. Furthermore, for $h\in\l B_0[z]_f$, we have $-\frac{1}{t^2}h(\frac{1}{t})\in\l B[t]_{f_1}$ and $\left(-\frac{1}{t^2}h(\frac{1}{t})\right)^c=-\frac{1}{t^2}h^c(t)$ for any $c\in\Hom_k(\l B,k)$. Thus $\left(-\frac{1}{t^2}A(\frac{1}{t})\right)^c=-\frac{1}{t^2}A(\frac{1}{t})^c$ and  $\del(y)=-\frac{1}{t^2}A(\frac{1}{t})^cy$ is regular singular, i.e., $\del(y)=A^cy$ is regular singular at $\infty$ for every $c\in\Hom_k(\l B,k)$. It follows that $\l B$ and $f$ have the required properties.
%	 
%	 we claim that the lemma holds with $\l B=\l B_1$ and $f=zf_0(z)f_2(z)f_1(z)\in\l B[z]$. Indeed, as 
%	 
%	 
%	  Enlarging $\l B_1$ if necessary, we can assume that $\l B_0\subseteq \l B_1$. Then,
%	 
%	  claim of the lemma then holds with $\l B=\l B_1$ and $f=f_0(z)f_1(z)$.
%	 
%	For the point $p=\infty$ we proceed similarly. % If $p=\infty$ is a regular (or regular singular) point for $\del(y)=Ay$
%	As $\del(y)=-\frac{1}{t^2}A(\frac{1}{t})y$ is regular singular, there exists, by Lemma~\ref{lem: specialize regsing at 0}, a finitely generated $k$-subalgebra $\l B_\infty$ of $k'$ and a monic polynomial $f_\infty\in\l B_\infty[t]$ such that $-\frac{1}{t^2}A(\frac{1}{t})\in \l B_\infty[t]_{f_\infty}^{n\times n}$ and $\del(y)=-\frac{1}{t^2}A(\frac{1}{t})^cy$ is regular singular for every $c\in\Hom_k(\l B_\infty,k)$. We may assume that $\l B_0\subseteq \l B_\infty$. Then $-\frac{1}{t^2}A(\frac{1}{t})^c=-\frac{1}{t^2}A^c(\frac{1}{t})$ and we see that $p=\infty$ is a regular singular point for $\del (y)=A^cy$ for all $c\in\Hom_k(\l B_\infty,k)$.
%	
%	Let $\l B$ be the $k$-subalgebra of $k'$ generated by $\l B_1,\ldots,\l B_m,\l B_\infty$ and set $f=f_0f_1\ldots f_mf_\infty\in\l B[z]$. We claim that $\l B$ and $f$ have the desired properties.
%	
%	
%	
%	Clearly $A\in\l B[z]_f^{n\times n}$. Let $c\in\Hom_k(\l B,k)$. As $\l B_0\subseteq \l B$ (so that $c$ restricts to a morphism $\l B_0\to k$), it follows from the first two paragraphs, that the poles of $\del(y)=Ay$ are contained in $S$. Similarly, as $c$ restricts to morphisms on $\l B_1,\ldots,\l B_m, \l B_\infty$, it follows from the above that $a_1,\ldots,a_m,\infty$ are regular singular points for $\del(y)=A^cy$.
%	As the poles of $\del(y)=A^cy$ are contained in $\{a_1,\ldots,a_m,\infty\}$, and all these points are regular singular, it follows that $\del(y)=A^cy$ is regular singular.	
%%	
%%	We claim that $\del(y)=A^cy$ is regular singular with singularities in $S$.
%	Note that $c\in\Hom_k(\l B,k)$ restricts to $c\in\Hom_k(\l B_i,k)$ for $k\in\{1,\ldots,m,\infty\}$. It follows from the above that $a_1^c,\ldots,a_r^c$ are regular points and $a_{r+1}^c,\ldots,a_m^c$ are regular singular points for $\del(y)=A^cy$. Moreover, $p=\infty$ is a regular (or regular singular) point for $\del(y)=A^cy$ if $p=\infty$ is a regular (or regular singular) point for $\del(y)=Ay$.
%	
%	
%	As $(z-a_1^c)\ldots(z-a_m^c)A^c=B^c\in  k[z]^{n\times n}$, the singularities of $\del(y)=A^cy$ are contained in $\{a_1^c,\ldots,a_m^c,\infty\}$. Because $a_1^c,\ldots,a_r^c$ are regular points for $\del(y)=A^cy$, in fact, the singularities of $\del(y)=A^cy$ are contained in $\{a_{r+1}^c,\ldots,a_m^c,\infty\}$. As all these points are regular singular for $\del(y)=A^c y$, we see that $\del(y)=A^c y$ is regular singular.
%	
%	By assumption, the singularities $a_{r+1},\ldots,a_m$ of $\del(y)=Ay$ are in $k$. So $a_i^c=a_i\in S$ for $i=r+1,\ldots,m$. If $\infty \in S$, this shows that the singularities of $\del(y)=A^cy$ are contained in $S$. If $\infty \notin S$, then $p=\infty$ is a regular point for $\del(y)=Ay$ and therefore $p=\infty$ is a regular point for $\del(y)=A^cy$. Thus, also in this case, the singularities of $\del(y)=A^cy$ are contained in $S$.
%	
%	Note that it is possible that a singularity of $\del(y)=Ay$ is a regular point for $\del(y)=A^cy$. This is no problem as we are only claiming that the singularities are contained in $S$.
%%	
%%	
%	
%	First, notice that the singularities of $A^c$ will be a subset of the singularities of $A$ : since they all lie in $\P^1(k)$, applying $c$ will leave them unchanged. Some singularities might disappear (say, if we have a coefficient of the form $a/z$ and $a$ is mapped to $0$) but that's not a problem because it just means we have less singularities to check.\\
%	We need to prove that $A^c$ is regular singular locally at every point, i.e. that $A^c$ is locally gauge equivalent to a matrix with a pole of order at most $1$ at every point. Since $A$ has only finitely many singularities, we prove that every singularity contributes a finite amount of generators to $\l B$ and a monic factor to $f$. We assume via a coordinate change that the singularity is at $0$.\\
%	\smallskip\noindent
%	Since $A$ is regular singular, Lemma \ref{lem:basis} ensures the existence of a $F \in \GL_{n}\big(\C(z)\big)$ and a matrix with regular coefficients $zB$ such that $F^{-1}(z\frac{d}{dz} - zA) F = z\frac{d}{dz} - zB$.\\
%	Moreover, there exists a monic polynomial $f_1 \in \C[z]$ such that $f_1 A$ and $f_1 F$ are in $\C[z]^{n \times n}$, ie $A$ and $F$ are in $\C[z]_{f_1}^{n \times n}$. We add all the coefficients of $f_1$, as well as those of the polynomial coefficients of $f_1A$ and $f_1F$ to $\l B_1$, so $A,F \in \l B_1[z]_{f_1}$\\
%	Next, we want $F$ to be invertible : we replace $f_1$ with $f_1\det(f_1F) \in \l B_1[z]$, such that :
%	\begin{align*}
%		F^{-1} &= \det(F)^{-1} \text{adj}(F) \\
%		&= f_1 \det(f_1F)^{-1} \text{adj}(f_1F)
%	\end{align*}
%	and by adding the inverse of the leading coefficient of this polynomial to $\l B_1$, we can assume it is monic (and will thus not specialize to the zero polynomial).\\
%	We construct the algebra $\l B$ by iterating this process on other singularities.\\
%	Applying the specialization $c$, we get that :
%	$$F^{c}\lp z\frac{d}{dz}-zA^c\rp \big(F^c\big)^{-1} = z\frac{d}{dz} - zB^c.$$
%	Since $zB^c$ has regular coefficients, the matrix $A^c$ is locally gauge equivalent to matrices with poles of order at most $1$, meaning it is regular singular.
\end{proof}

\end{document}
